% ============================================================================
\section{Does Supervision Density Change Update Sparsity?}
\label{sec:density}
% ============================================================================

Our third study asks whether vocabulary-level supervision changes how
strongly parameter updates concentrate on a small fraction of model parameters.
Sampled-token PG uses feedback for the generated token, full-vocabulary
reverse KL includes every vocabulary alternative, and teacher top-$k$
provides a practical intermediate choice
(Section~\ref{sec:setting}). This concerns vocabulary coverage at each
prefix, rather than the number of supervised response positions.
Sampled-token feedback does not imply a sparse parameter gradient:
the sampled token's log probability depends on all logits through the
softmax normalizer. Concentrated OPD parameter changes have also been
observed with dense feedback \citep{yu2026denseupdates}. We therefore
test for a difference without assuming that PG produces sparser updates
or better capability integration.

\subsection{Matched single-teacher and joint experiments}

We select one representative single-teacher setting before inspecting
results and compare PG with corrected teacher top-$64$ in that setting
and in joint MOPD (Table~\ref{tab:density_controls}). These four online
configurations use the losses in Equations~\ref{eq:mopd_pg}
and~\ref{eq:mopd_topk}; the joint top-$64$ run also anchors
Section~\ref{sec:normalization}. All use domain-response averaging and
match initialization, teachers, optimizer, loss masks, response cap,
prompt schedule, and update budget within each setting. The top-$64$
loss uses probabilities normalized over the full vocabulary, without
renormalizing within the selected tokens; we record any clipping.
The shared model and training setup appears in
Section~\ref{sec:normalization} and Appendix~\ref{app:setup}.

\begin{table}[!htbp]
\centering
\small
\begin{tabular}{lll}
\toprule
Setting & Online losses & Common-prefix reference \\
\midrule
Single teacher & PG, teacher top-$64$ & Full-vocabulary reverse KL \\
Joint MOPD & PG, teacher top-$64$ & Full-vocabulary reverse KL \\
\bottomrule
\end{tabular}
\caption{Supervision-density experiments. Four online configurations
share full-vocabulary local measurements at matched checkpoints.}
\label{tab:density_controls}
\end{table}

At the initial, early, intermediate, and final checkpoints, we evaluate
all three losses on a small shared bank of student-generated prefixes.
Each comparison holds the student, prefix weights, and optimizer state
fixed; local PG uses fresh token samples at those prefixes.
Full-vocabulary reverse-KL proposed steps provide the local reference in
both settings, with raw gradients as a diagnostic; these local measurements
cannot establish the cumulative changes of a full-vocabulary online run,
which is conditional on measured resource feasibility. Online PG and
top-$64$ branches generate their own on-policy responses, so their
trajectories compare practical loss choices rather than isolating
vocabulary coverage as the only causal difference.

\subsection{Update concentration and learning outcomes}

Our primary measurements are actual online steps and cumulative parameter
changes from initialization, reported separately. Following
Section~\ref{sec:subnetworks}, we report threshold curves, norms, and the
fraction of squared change carried by the largest coordinates. To measure
effective cumulative changes, we compare stored BF16 checkpoints with
initialization at absolute tolerance $10^{-5}$, following
\citet{mukherjee2025subnetworks}. FP32 master-weight differences separately
characterize accumulated optimizer motion. Local comparisons use proposed steps from
copies of the same optimizer state; repeated batches on the common prefix
bank quantify sampling variability. Raw gradients before clipping provide
secondary concentration measurements. Their sparsity need not equal update
sparsity because momentum, clipping, weight decay, rounding, and cancellation
across steps can change which parameters move.

For the joint runs, we also compare which parameters the teachers' proposed
steps select, with raw-gradient overlap as a diagnostic, using
Section~\ref{sec:subnetworks}'s overlap measurement. Held-out per-domain
capability, the macro-average,
and worst-domain change accompany the sparsity results. Curves show
generated-token exposure and GPU hours separately, with response counts,
updates, and teacher-scoring and diagnostic overhead recorded.

\paragraph{Preliminary observation: drift and effective BF16 changes.}
In seed-$42$ top-$64$ runs with domain-response (DR) averaging, logged
student mass inside the teacher's top-$64$ falls from $86.2\%$ to $37.0\%$
for the single teacher and from $81.4\%$ to $51.7\%$ for joint training
(rollout batches $1$--$50$ versus $200$--$249$). Joint domain-token and
global-token means instead rise to $93.7\%$ and $94.9\%$, respectively;
each run uses its own on-policy prefixes and configured averaging.
Under the BF16 criterion above, single-teacher top-$64$ changes
$5.33\%$ of parameters at update $100$, versus $2.55\%$ for PG.
Their BF16 cumulative-change $L_2$ norms are $0.201$ and $0.125$,
respectively.
At update $250$, joint top-$64$ DR changes $9.92\%$, versus
$5.90\%$, $11.72\%$, and $12.01\%$ for PG, domain-token, and global-token;
BF16 cumulative-change $L_2$ norms are $0.425$, $0.269$, $0.537$, and
$0.555$, respectively. Thus top-$64$ changes more effective BF16 weights
than PG. The stronger distribution drift under DR is consistent with an
interaction between truncation and response averaging, but does not
coincide with the largest BF16 parameter changes among the averaging
rules. The mechanism remains to be established.
% Probability windows and FP32 diagnostics: experiments/topk_drift_sparsity_20260906.json
% Full BF16 weight scan using the cited paper's criterion:
% experiments/bf16_effective_changes_20260906.json

\subsection{Interpretation and decision rule}

At a fixed prefix, fresh unclipped sampled-token PG with a detached
log-ratio coefficient has the same expected gradient as full-vocabulary
reverse KL \citep{yu2026denseupdates,oh2026vopd}. This identity holds with
fixed prefix weights; weighting an original rollout token by its eventual
response length need not preserve it. Clipping and top-$k$ truncation
also change the comparison, and equality in expectation does not equate
sampled gradients, nonlinear optimizer steps, or training trajectories.

We report paired-seed effect estimates and confidence intervals against
a material difference margin specified before examining results.
One-seed runs remain exploratory; prompt resampling measures evaluation
uncertainty, not training-seed variation. Intervals entirely within the
margin support no material sparsity difference in the tested regime;
intervals spanning negligible and material effects are inconclusive.
Only a stable, meaningful gap motivates
the optional sampling-variance and truncation controls in
Appendix~\ref{app:optional_loss}.

\missing{PG/top-$64$ optimizer steps, cumulative changes, and capability;
full-vocabulary local references and raw-gradient diagnostics.}
